\documentclass[11pt]{article}

% ---------- packages ----------
\usepackage[utf8]{inputenc}
\usepackage[T1]{fontenc}
\IfFileExists{lmodern.sty}{\usepackage{lmodern}}{}
\usepackage[margin=1in]{geometry}
\usepackage{amsmath,amssymb}
\usepackage{graphicx}
\usepackage{booktabs}
\usepackage{array}
\usepackage{caption}
\usepackage{authblk}
\usepackage[hidelinks]{hyperref}

% Show a placeholder box if a figure file is not present, so the document
% always compiles. Drop the real PDF (e.g. fullplanck_w0wa.pdf) next to this
% .tex file and it will be used automatically.
\newcommand{\figorbox}[3]{%
  \IfFileExists{#1}{\includegraphics[width=#2]{#1}}%
  {\fbox{\parbox[c][4.5cm][c]{#2}{\centering\itshape [Insert figure here]\\[4pt]\normalfont\small #3}}}%
}

\captionsetup{font=small,labelfont=bf}
\setlength{\parskip}{4pt}

% ---------- title ----------
\title{\bfseries A Low-$z$ / Mid-$z$ Nuisance Mode in the Union3 Compressed
Supernova Data: Implications for the DESI DR1 Dynamic Dark Energy Preference}

\author{Jaanus Hass}
\affil{Independent researcher, Saru, Estonia \\ \texttt{jaanush@hotmail.com}}
\date{May 2026 \quad(full-likelihood cross-check added June 2026)}

\begin{document}
\maketitle

\begin{center}
\small\itshape
Analysis performed May--June 2026 using AI-assisted coding and checking.
All results are reproducible from public data.
\end{center}

% ====================================================================
\begin{abstract}
\noindent
We test whether the Union3 supernova contribution to the DESI DR1 CPL dynamical
dark-energy preference is robust to a single redshift-dependent nuisance direction.
Using the public Union3 compressed 22-node precision matrix and individual
light-curve fit files, we find that a Gaussian low-$z$ / mid-$z$ nuisance template
$T(z)=G(z;0.09,0.06)-G(z;0.775,0.10)$ improves the fixed-$\Omega_m=0.2975$ fit by
$\Delta\chi^2=17.12$ (AIC $+15.12$), compared with $\Delta\chi^2\approx7.53$
(AIC $\approx+3.53$) for CPL. After including this nuisance, CPL adds only
$\Delta\chi^2\approx0.87$ and is AIC-disfavored.
A public-data Cobaya cross-check using DESI DR1 BAO, a compressed CMB prior on
$\Omega_m h^2$, and the Union3 compressed likelihood gives the same qualitative
result: CPL improves $\chi^2$ by $8.29$ without the nuisance mode, but only by
$0.52$ after the nuisance mode is included, with the CPL+nuisance best-fit solution shifting
close to $\Lambda$CDM ($w_0=-0.999$, $w_a=-0.192$).
\textbf{A new full-likelihood cross-check}---Planck NPIPE CamSpec TTTEEE plus
low-$\ell$ TT/EE, DESI DR1 BAO, and the Union3 likelihood with the nuisance amplitude
sampled directly (four MCMC chains; the means converged to $R-1=0.047$, the bounds criterion more weakly, $R-1\approx0.15$)---confirms this: without the nuisance,
CPL is preferred over $\Lambda$CDM by AIC gain $+8.42$ ($w_0=-0.677\pm0.113$,
$w_a=-1.136\pm0.438$); once the nuisance is marginalized the CPL preference is removed
(AIC gain $-0.38$ relative to $\Lambda$CDM$+$nuisance), with the posterior shifting to
$w_0=-0.933\pm0.129$, $w_a=-0.501\pm0.400$ and a recovered amplitude
$a=0.029\pm0.010$ consistent with the compressed-data value.
Additional robustness tests (81 template variants; a 200-template random null test;
leave-one-region-out; cumulative low-$z$ node removal) strengthen the diagnostic.
The dynamical-dark-energy preference is weakened from $\sim2.4\sigma$ to $\sim1.6\sigma$
but not entirely eliminated, consistent with the partial weakening already seen in the
official Union3.1 joint analyses. The definitive test remains to rerun the official
DESI$+$CMB$+$Union3 $w_0$--$w_a$ likelihood with an added redshift-dependent nuisance
parameter.
\end{abstract}

% ====================================================================
\section{Introduction}

The DESI 2024 Data Release 1 BAO analysis, combined with CMB and Type~Ia supernova
data, reports a preference for dynamical dark energy described by the CPL equation of
state $w(a)=w_0+w_a(1-a)$, with a combined significance of approximately
$3.4$--$3.5\sigma$ for DESI$+$CMB$+$Union3 \cite{desi2024}. The Union3 compilation
\cite{rubin2025} contains 2,087 Type~Ia supernovae from 24 surveys spanning
approximately 7 billion years of cosmic history, compressed into a 22-node
distance-modulus spline with an associated precision matrix. Individual SN surveys
cover limited redshift ranges, making multi-survey compilations necessary but
potentially introducing redshift-dependent systematics at survey boundaries. Related
work by Efstathiou \cite{efstathiou2025} identified a systematic offset between low-
and high-redshift supernovae in Pantheon$+$ and DES5Y, and noted that Union3
individual-SN data were not publicly available for direct analysis. Subsequent analyses have sharpened this question: Huang et al.\ \cite{huang2025} used the intercept of the SN magnitude--redshift relation to argue that correcting a low-$z$ discrepancy in DES5Y reduces the CPL preference to $0.5$--$1.5\sigma$, and Ormondroyd et al.\ \cite{ormondroyd2026} showed via Bayesian model selection that a low-$z$ magnitude offset is favoured over a dynamical reconstruction for DES5Y$+$DESI; most recently Mirpoorian et al.\ \cite{mirpoorian2026} fit a look-back-time-scaled magnitude nuisance jointly with $w_0w_a$CDM on Pantheon$+$, finding negligible impact for uncalibrated SNe. These interpretations remain contested, with Vincenzi et al.\ \cite{vincenzi2025} attributing much of the offset to Malmquist-bias and intrinsic-scatter modelling and the recalibrated DES ``Dovekie'' analysis \cite{popovic2025} continuing to prefer evolving dark energy. Almost all of this work targets DES5Y or Pantheon$+$, whose per-object data are public, which is what makes the Union3-specific test below distinct. Cort\^es \&
Liddle \cite{cortes2025}, addressing DESI DR2, further showed that properly accounting
for the substantial overlap between primary SN compilations reduces the combined
$\Lambda$CDM exclusion significance to $3.1\sigma$---identical to the DESI BAO$+$CMB
result without any supernova data; their methodological argument about compilation
overlap is relevant to DR1-style compilation comparisons as well. Independent
forecasting work by Truong et al.\ \cite{truong2026}, evaluated at the same
DESI$+$CMB$+$Union3 best fit, shows that the $w_0$--$w_a$ discriminating power of
SNe~Ia is concentrated at $z\lesssim0.6$, with low-redshift supernovae primarily
constraining $w_0$ and a medium-redshift anchor required to make them informative;
a coherent low-$z$/mid-$z$ distance distortion is therefore precisely the kind of
systematic direction that can mimic or absorb the dynamical-dark-energy signal,
which motivates the template form adopted here. The present work
performs, to the best of the author's knowledge, the first compressed-data robustness
test of this kind specifically for Union3, using the publicly available
\texttt{lcfit\_Union3.tar.gz} and compressed precision matrix from the
\texttt{rubind/union3\_release} repository.

% ====================================================================
\section{Data and Method}

We use the public files \texttt{mu\_mat\_union3\_cosmo=2\_mu.fits} and
\texttt{lcfit\_Union3.tar.gz} from the \texttt{rubind/union3\_release} GitHub
repository \cite{data}. The FITS matrix structure is: \texttt{mat[0,1:]} $=z$ nodes (22 values),
\texttt{mat[1:,0]} $=\mu$ values, \texttt{mat[1:,1:]} $=$ precision matrix $P$
($22\times22$). The magnitude offset $M$ is analytically marginalized in all fits via
GLS; the nuisance amplitude $a$ is also analytically marginalized when included. For
node-removal tests, the covariance matrix $C=P^{-1}$ is subselected and re-inverted
correctly. The primary nuisance template is
\begin{equation}
T(z)=G(z;0.09,0.06)-G(z;0.775,0.10),
\end{equation}
standardized to zero mean and unit RMS, where $G(z;c,s)=\exp[-\tfrac12((z-c)/s)^2]$.
Template centers were placed at known survey-composition transitions (not optimized
over the data): the low-$z$ anchor region dominated by SDSS/Foundation/Pan-STARRS
($z\approx0.09$) and the mid-$z$ transition into SNLS/ESSENCE/DES3\_Deep
($z\approx0.775$). Quality-cut pass fractions differ between these zones: $0.718$ at
$z\approx0.10$ versus $0.952$ at $z\approx0.75$.

For the public joint cross-check we use the Cobaya DESI DR1 BAO likelihood, a
compressed CMB prior $\Omega_m h^2=0.1430\pm0.0020$ (Planck 2018 \cite{planck2018}; deliberately
broadened from the Planck best-fit $\sigma=0.0011$ as a conservative choice), and the
Union3 compressed likelihood. Sampled parameters: $H_0$ (flat prior 55--80
km/s/Mpc), $\Omega_m$ (flat 0.1--0.6), $w_0$ (flat $-3$ to $1$), $w_a$ (flat $-3$ to
$2$), and nuisance amplitude $a$ (flat $-0.5$ to $0.5$ mag). Flat geometry and standard
neutrino content ($\Sigma m_\nu=0.06$ eV) assumed; sound horizon fixed at
$r_d=147.09$ Mpc. Within-redshift $D_M$--$D_H$ BAO correlations are included;
between-redshift-bin correlations are set to zero. We note that Kim \cite{kim2024}
showed that the Union3 compressed product distributes a posterior under a spline-node
prior rather than a raw likelihood, which explains why proxy baseline CPL best-fits
differ slightly from the official DESI result; this does not affect the relative CPL
vs CPL$+$nuisance comparison.

% ====================================================================
\section{Union3 Compressed-Data Results}

Table~\ref{tab:protocol} summarises the protocol verification. The core fixed-$\Omega_m$
diagnostics reproduce to rounding precision. A low-minus-mid step template also gives a
strong improvement, confirming that the effect is not specific to the Gaussian
functional form.

\begin{table}[h]
\centering
\caption{Protocol verification summary.}
\label{tab:protocol}
\small\setlength{\tabcolsep}{4pt}
\begin{tabular}{lccl}
\toprule
Test & Claimed & Verified & Status \\
\midrule
$\Omega_m$ best-fit ($\Lambda$CDM free) & 0.356 & 0.3559 & Reproduced \\
$\chi^2$ at best-fit & 23.96 & 23.958 & Reproduced \\
DESI-like $\Omega_m$ penalty & 5.15 & 5.152 & Reproduced \\
CPL $\Delta\chi^2$ / AIC & 7.53 / $+3.53$ & 7.536 / $+3.536$ & Reproduced \\
Low-mid Gaussian $\Delta\chi^2$ / AIC & 17.12 / $+15.12$ & 17.120 / $+15.120$ & Reproduced \\
Color $\Delta\chi^2$ / AIC & 6.52 / $+4.52$ & 2.947 / $+0.947$ & Weighting-dependent \\
CPL after low-mid / AIC & 0.87 / $-3.13$ & 0.870 / $-3.130$ & Reproduced \\
Low-mid prior threshold $\sigma_a\ge0.005$ mag & 0.005 mag & RMS & Reproduced \\
\bottomrule
\end{tabular}
\end{table}

% ====================================================================
\section{Joint DESI DR1 BAO + CMB + Union3 Public Cross-Check (compressed CMB prior)}

We ran a public-data Cobaya cross-check using the DESI DR1 BAO likelihood, a compressed
CMB prior, and the Union3 compressed likelihood. This is stronger than a simple proxy
grid scan, but it is still not the official DESI$+$CMB$+$Union3 full likelihood because
the CMB information is compressed and the Union3 nuisance parameter is added externally.

\begin{table}[h]
\centering
\caption{Public-data joint model comparison (compressed CMB prior). AIC gain
$=\Delta\chi^2-2k$, where $k$ is the number of additional parameters vs $\Lambda$CDM.
The highlighted row shows that CPL adds only $\Delta\chi^2=0.52$ once the nuisance
direction is included.}
\label{tab:proxy}
\begin{tabular}{lcccccc}
\toprule
Model & $H_0$ & $w_0$ & $w_a$ & $a_\mathrm{nuis}$ & best-sample $\chi^2$ & $\Delta\chi^2$ / AIC gain \\
\midrule
$\Lambda$CDM            & 68.77 & ---     & ---     & ---    & 41.24 & --- \\
CPL                     & 66.64 & $-0.688$ & $-1.070$ & ---   & 32.95 & $+8.29$ / $+4.29$ \\
$\Lambda$CDM $+$ nuis.  & 69.09 & ---     & ---     & 0.0279 & 25.04 & $+16.20$ / $+14.20$ \\
CPL $+$ nuisance        & 69.79 & $-0.999$ & $-0.192$ & 0.0306 & 24.51 & $+16.73$ / $+10.73$ \\
\midrule
CPL vs $\Lambda$CDM$+$nuis. & & & & & & $+0.52$ / $-3.48$ \\
\bottomrule
\end{tabular}
\end{table}

Without the nuisance mode, CPL improves $\chi^2$ by $8.29$ and shifts to $w_0=-0.688$,
$w_a=-1.070$, consistent with published DESI DR1$+$CMB$+$Union3 results. After including
the nuisance, CPL improves $\chi^2$ by only $0.52$ relative to $\Lambda$CDM$+$nuisance
(AIC $-3.48$), while the CPL$+$nuisance best-fit solution shifts to $w_0=-0.999$, $w_a=-0.192$---close
to $\Lambda$CDM in this proxy analysis.

% ====================================================================
\section{Full-Likelihood Cross-Check (Planck NPIPE CamSpec + DESI DR1 BAO + Union3)}
\label{sec:fullplanck}

We repeated the four-model comparison with the full likelihood, replacing the
compressed CMB prior of Section~4 with Planck NPIPE CamSpec TTTEEE plus low-$\ell$
TT/EE, and adding the redshift-dependent nuisance parameter directly inside the Union3
likelihood (amplitude $a$). Four MCMC chains were run per model; the Gelman--Rubin
convergence on the parameter means reached $R-1=0.047$, while the more stringent
bounds (credible-interval) criterion was weaker, $R-1\approx0.15$. The full-likelihood
comparison should therefore be regarded as a public cross-check rather than
collaboration-grade parameter estimation. This realizes the falsification test of
Section~10 within the public pipeline, with the nuisance mode marginalized rather
than added externally.

\begin{table}[h]
\centering
\caption{Full-likelihood model comparison (Planck NPIPE CamSpec TTTEEE $+$ low-$\ell$
TT/EE $+$ DESI DR1 BAO $+$ Union3). Four MCMC chains per model (means $R-1=0.047$,
bounds $R-1\approx0.15$). Quoted $\chi^2$ are the best values among the MCMC samples
(no separate minimizer was run); the present table is therefore an MCMC best-sample
cross-check, and a separate minimizer run would be required for final best-fit
publication numbers. AIC gain $=\Delta\chi^2-2k$ relative to
$\Lambda$CDM; positive $=$ favoured. The highlighted comparison shows that once the
nuisance direction is marginalized, the CPL parameters ($w_0,w_a$) no longer improve AIC.}
\label{tab:fullplanck}
\small\setlength{\tabcolsep}{3pt}
\begin{tabular}{lcccccc}
\toprule
Model & $H_0$ & $w_0$ & $w_a$ & $a_\mathrm{nuis}$ & best-sample $\chi^2$ & $\Delta\chi^2$ / AIC gain \\
\midrule
$\Lambda$CDM            & 67.8 & ---              & ---              & ---                & 11007.56 & --- \\
CPL                     & 66.4 & $-0.677\pm0.113$ & $-1.136\pm0.438$ & ---                & 10995.14 & $+12.42$ / $+8.42$ \\
$\Lambda$CDM $+$ nuis.  & 68.0 & ---              & ---              & $0.027\pm0.010$    & 10992.23 & $+15.33$ / $+13.33$ \\
CPL $+$ nuisance        & 69.3 & $-0.933\pm0.129$ & $-0.501\pm0.400$ & $0.029\pm0.010$    & 10988.61 & $+18.95$ / $+12.95$ \\
\midrule
\multicolumn{6}{l}{CPL $+$ nuisance \emph{vs} $\Lambda$CDM $+$ nuisance}
                        & $+3.62$ / $\mathbf{-0.38}$ \\
\bottomrule
\end{tabular}
\end{table}

Without the nuisance mode, CPL is strongly preferred over $\Lambda$CDM
($\Delta\chi^2=12.42$, AIC gain $+8.42$), with $w_0=-0.677\pm0.113$,
$w_a=-1.136\pm0.438$---a posterior mean lying $2.4\sigma$ from $\Lambda$CDM in the
$w_0$--$w_a$ plane (Gaussian-equivalent of the 2D Mahalanobis distance $r=2.86$),%
\footnote{The equivalent one-dimensional significance is the two-sided Gaussian tail
enclosing the same probability as the 2D Mahalanobis ellipse,
$P=1-\exp(-r^2/2)$; i.e.\ $\mathrm{erf}(\sigma/\sqrt2)=1-\exp(-r^2/2)$. This maps
$r=2.86\to2.4\sigma$ and $r=2.11\to1.6\sigma$.}
consistent with the published DESI DR1$+$CMB$+$SN preference. Including the nuisance
parameter yields $a=0.029\pm0.010$ ($\sim2.8\sigma$ from zero, consistent with the
compressed-data amplitude) and shifts the posterior to $w_0=-0.933\pm0.129$,
$w_a=-0.501\pm0.400$, now $1.6\sigma$ from $\Lambda$CDM ($r=2.11$).

The decisive comparison is CPL$+$nuisance versus $\Lambda$CDM$+$nuisance. Once the
nuisance direction is marginalized, adding the CPL parameters improves the best-fit
$\chi^2$ by only $3.62$ for two extra parameters, corresponding to an AIC gain of
$-0.38$: the strong AIC preference for dynamical dark energy ($+8.42$ relative to
$\Lambda$CDM without the nuisance) is removed. The full-likelihood result therefore
confirms the compressed-data analysis: the Union3 low-$z$/mid-$z$ nuisance direction is
favoured by the data and absorbs most of the dynamical-dark-energy preference. As in the
compressed proxy, the effect weakens but does not entirely eliminate the preference (the
shift is from $\sim2.4\sigma$ to $\sim1.6\sigma$, and the residual CPL AIC penalty is
marginal), consistent with the partial weakening already seen in the official Union3.1
joint analyses (Section~6).

\begin{figure}[h]
\centering
\figorbox{fullplanck_w0wa_desi.pdf}{0.75\textwidth}{w0-wa: DESI official (black), this work CPL (blue), CPL+nuisance (green), LCDM star.}
\caption{Full-likelihood joint cross-check in the $w_0$--$w_a$ plane. Black contours:
DESI DR1 official $w_0w_a$CDM posterior (DESI$+$CMB$+$Union3, \texttt{plik} TT,TE,EE).
Blue: this work's CPL (NPIPE CamSpec); green: this work's CPL$+$nuisance. Red star marks
$\Lambda$CDM ($w_0=-1$, $w_a=0$). This work's CPL reproduces the official DESI contour
($w_0=-0.68$, $w_a=-1.14$ vs DESI $w_0=-0.64$, $w_a=-1.30$), the small offset being
consistent with the different CMB likelihood choice; marginalizing the Union3 nuisance
direction (green) shifts the posterior toward $\Lambda$CDM ($w_0=-0.93$, $w_a=-0.50$),
from $\sim2.4\sigma$ to $\sim1.6\sigma$.}
\label{fig:fullplanck}
\end{figure}

% ====================================================================
\section{Three-Version Progression: Union3 to Union3.1}

Hoyt et al.\ \cite{hoyt2026} published host-galaxy mass re-measurements for
approximately 2,000 Union3 supernovae, finding that low-$z$ ($<0.10$) host-galaxy
stellar masses were systematically overestimated, reducing a previously reported
$>0.03$ mag low-$z$ distance offset to $\approx0.01$ mag. The updated cosmological
framework (UNITY1.8) was presented by Rubin et al.\ \cite{rubin2026}. Applying the same
nuisance template to all three public compressed matrices shows the nuisance amplitude
decreasing monotonically, while the CPL AIC gain becomes small by AIC comparison in the
Union3.1 versions. The nuisance amplitude decrease of $0.0056$ mag
($0.0279\to0.0223$) is of the same order of magnitude as the relative low-$z$ distance
shift of $0.010$--$0.013$ mag applied in the Union3.1 corrections. This correspondence
was not built into the template---the amplitude emerged from fitting the compressed
data, and is consistent with the template tracing a low-$z$ systematic direction that
is reduced, but not removed, by the Union3.1/UNITY processing. This monotonic decrease
also bears directly on the central interpretive question of whether the template
absorbs a genuine systematic or, instead, real cosmological signal: the underlying
$w_0$--$w_a$ cosmology is identical across the three compressed matrices, so a template
tracing a purely cosmological signal would not naturally be expected to decrease in step with documented low-$z$ host-mass corrections. Instead
the amplitude falls in step with the documented low-$z$ host-mass correction of Hoyt
et al.\ \cite{hoyt2026}---the behaviour expected if the template is tracing that
corrected systematic rather than the dynamical-dark-energy signal. Importantly, this
comparison is a heuristic compressed-data diagnostic only and should not be conflated
with the updated official joint analyses: Hoyt et al.\ \cite{hoyt2026} report a
$\approx3.4\sigma$ deviation from $\Lambda$CDM in the official BAO$+$CMB$+$Union3.1
combination, and Rubin et al.\ \cite{rubin2026} report $\approx2.6\sigma$ tension with
UNITY1.8. The official joint preference is thus weakened but not eliminated by the
Union3.1 corrections.

\begin{table}[h]
\centering
\caption{Three-version comparison. [a] CPL $\Delta\chi^2$ values marked $\approx$ are
grid-scan approximations with $\approx0.02$ uncertainty; the scipy-optimize value for
Union3 original is 7.536 (Table~\ref{tab:protocol}). Nuisance amplitude decreases
monotonically across all three versions.}
\label{tab:threeversion}
\begin{tabular}{lccc}
\toprule
Metric & Union3 original & Union3.1 UNITY1.7 & Union3.1 UNITY1.8 \\
\midrule
$\Omega_m$ best-fit & 0.3559 & 0.3424 & 0.3340 \\
DESI penalty $\Delta\chi^2$ & 5.152 & 3.256 & 2.347 \\
CPL $\Delta\chi^2$ [a] & $\approx7.52$ & $\approx4.79$ & $\approx4.74$ \\
CPL AIC gain [a] & $\approx+3.52$ & $\approx+0.79$ & $\approx+0.74$ \\
Low-mid Gaussian $\Delta\chi^2$ & 17.120 & 12.474 & 11.958 \\
Low-mid Gaussian AIC gain & $+15.12$ & $+10.47$ & $+9.96$ \\
Nuisance amplitude (mag) & 0.0279 & 0.0233 & 0.0223 \\
CPL after nuisance AIC & $-3.13$ & $-2.80$ & $-2.04$ \\
Low-$z$ shift vs Union3 & --- & $-0.010$ & $-0.013$ \\
\bottomrule
\end{tabular}
\end{table}

% ====================================================================
\section{Additional Robustness Tests}

The robustness tests in this section use the same fixed-$\Omega_m=0.2975$ protocol as
Table~\ref{tab:protocol}; their $\Delta\chi^2$ values are directly comparable to the
Table~\ref{tab:protocol} results and are reproducible from the GitHub repository CSV
outputs. CPL$+$nuisance $\Delta\chi^2$ values in reduced-node fits can be
optimizer-sensitive and should not be interpreted as precise parameter estimates.

\subsection{Template Robustness Grid}

To test whether the result depends on the exact Gaussian centers and widths, we scanned
81 nearby low-$z$ / mid-$z$ template variants. The nuisance improvement remains positive
in all 81 cases. The conclusion does not depend on the primary template sitting near
the top of this grid (original $17.12$ vs grid maximum $17.45$): even the grid
\emph{median} ($13.18$) comfortably exceeds the CPL improvement ($7.54$), so any
representative low-$z$/mid-$z$ template in this family out-performs CPL.

\begin{table}[h]
\centering
\caption{Template robustness grid summary.}
\label{tab:grid}
\begin{tabular}{lc}
\toprule
Quantity & Value \\
\midrule
Number of templates & 81 \\
$\Lambda$CDM $\chi^2$ (fixed $\Omega_m$) & 29.110 \\
CPL $\chi^2$ (fixed $\Omega_m$) & 21.574 \\
CPL $\Delta\chi^2$ vs $\Lambda$CDM & 7.536 \\
Original nuisance $\Delta\chi^2$ & 17.120 \\
Grid median nuisance $\Delta\chi^2$ & 13.176 \\
Grid min nuisance $\Delta\chi^2$ & 9.666 \\
Grid max nuisance $\Delta\chi^2$ & 17.454 \\
Median CPL $\Delta\chi^2$ after nuisance & 0.413 \\
\bottomrule
\end{tabular}
\end{table}

\subsection{Random Smooth-Template Null Test}

A random smooth-template null test compares the original low-$z$ / mid-$z$ template with
200 random smooth redshift templates generated as follows: each random template is a
difference of two Gaussians with centers drawn uniformly from $z\in[0.05,0.40]$ and
$z\in[0.40,1.50]$ respectively, and widths drawn uniformly from $[0.04,0.20]$, with
random sign; each template is then standardized to zero mean and unit RMS on the 22
nodes before fitting (random seed 42 for reproducibility). No random template reached
the original template improvement; conservatively, the empirical tail probability is
$p\approx1/201$ (add-one estimator).

\begin{table}[h]
\centering
\caption{Random smooth-template null test (fixed-$\Omega_m=0.2975$ protocol). The
original template $\Delta\chi^2=17.120$; the maximum random template reached only
12.433. The 0/200 exceedance result is a conservative ranking diagnostic rather than
proof of template uniqueness.}
\label{tab:null}
\begin{tabular}{lc}
\toprule
Quantity & Value \\
\midrule
Number of random templates & 200 \\
Original template $\Delta\chi^2$ (fixed $\Omega_m$) & 17.120 \\
Random median $\Delta\chi^2$ & 1.357 \\
Random 95th percentile $\Delta\chi^2$ & 8.997 \\
Random max $\Delta\chi^2$ & 12.433 \\
Empirical tail probability & 0/200; add-one $p\approx1/201$ ($\approx0.005$) \\
\bottomrule
\end{tabular}
\end{table}

A sharper version of this test isolates the structure the template captures
\emph{beyond} CPL. For each template we compute the improvement it adds once CPL
is already in the model,
$\Delta\chi^2_{\text{after CPL}}=\chi^2(\text{CPL})-\chi^2(\text{CPL}+\text{template})$,
which removes the part of the residual that CPL can itself absorb (cf.
Section~\ref{sec:overlap}). Applied to the same 200-template ensemble, the
original low-$z$/mid-$z$ template adds $\Delta\chi^2=10.45$ after CPL, exceeding
\emph{every} random smooth template (maximum $8.99$, median $1.33$; add-one
$p\approx1/201$; Table~\ref{tab:nullaftercpl}). The structure the template
captures beyond CPL is therefore not generic overfitting of a flexible direction:
no random smooth template orthogonal-enough to CPL reaches it. As in the standalone
test this is a conservative ranking diagnostic rather than a proof of uniqueness.

\begin{table}[h]
\centering
\caption{Beyond-CPL (orthogonalized) null test: improvement added \emph{after} CPL,
$\Delta\chi^2_{\text{after CPL}}$, for the original template versus 200 random smooth
templates (same recipe and fixed-$\Omega_m=0.2975$ protocol as Table~\ref{tab:null}).
Removing the CPL-absorbable part lowers the random ceiling from $12.4$ to $9.0$, yet
the original template still exceeds all 200.}
\label{tab:nullaftercpl}
\begin{tabular}{lc}
\toprule
Quantity & Value \\
\midrule
Original template $\Delta\chi^2_{\text{after CPL}}$ & 10.45 \\
Random median $\Delta\chi^2_{\text{after CPL}}$ & 1.33 \\
Random 95th percentile $\Delta\chi^2_{\text{after CPL}}$ & 5.42 \\
Random max $\Delta\chi^2_{\text{after CPL}}$ & 8.99 \\
Empirical tail probability & 0/200; add-one $p\approx1/201$ ($\approx0.005$) \\
\bottomrule
\end{tabular}
\end{table}

\subsection{Leave-One-Redshift-Region-Out Test}

We removed broad redshift intervals one at a time and repeated the four-model
compressed-data fit. The nuisance improvement remains substantial in every cut, while
CPL adds little after nuisance is included. The largest weakening occurs when
$z=0.10$--$0.30$ is removed, localizing part of the signal to that low-redshift
transition region.

\begin{table}[h]
\centering
\caption{Leave-one-redshift-region-out robustness test. Only $\Delta\chi^2$ differences
are interpreted; cosmological parameter values in reduced-node fits are diagnostic and
may be weakly constrained.}
\label{tab:loo}
\begin{tabular}{lccccc}
\toprule
Case & Removed $z$ & $N$ kept & $\Delta\chi^2$ CPL & $\Delta\chi^2$ nuisance & CPL after nuisance \\
\midrule
Full sample   & none        & 22 & 7.54  & 17.12 & 0.87 \\
Very low $z$  & 0.00--0.10  & 21 & 7.58  & 16.76 & 0.82 \\
Low $z$       & 0.10--0.30  & 18 & 2.80  & 7.65  & 0.48 \\
Mid $z$       & 0.30--0.70  & 14 & 13.26 & 17.57 & 0.30 \\
High-mid $z$  & 0.70--1.00  & 17 & 7.19  & 12.10 & 0.46 \\
High $z$      & 1.00$+$     & 18 & 8.05  & 17.51 & 1.38 \\
\bottomrule
\end{tabular}
\end{table}

\subsection{Cumulative Low-$z$ Node Removal}

As a complement to the broad leave-one-region-out test, we progressively removed
low-redshift nodes below successively higher thresholds and re-evaluated CPL on the
remaining nodes. This directly tests the concern that the nuisance template may
effectively remove low-$z$ supernova information. Removing only the single lowest node
($z<0.08$ or $<0.10$) leaves CPL essentially unchanged ($\Delta\chi^2\approx7.58$). The
CPL preference falls sharply only when nodes up to $z<0.12$ are removed
($\Delta\chi^2=2.29$, AIC $-1.71$) and is nearly eliminated by $z<0.20$
($\Delta\chi^2=0.87$, AIC $-3.13$). This localizes the signal to the
$z=0.10$--$0.20$ transition region, consistent with the low-$z$ / mid-$z$ nuisance
template and not attributable to a single outlier node.

\begin{table}[h]
\centering
\caption{CPL preference after cumulative removal of low-redshift nodes. The sharp drop
between $z<0.10$ and $z<0.12$ localizes the signal to the $z=0.10$--$0.20$
survey-transition region.}
\label{tab:cumulative}
\begin{tabular}{lccc}
\toprule
Removed nodes & $N$ kept & CPL $\Delta\chi^2$ & CPL AIC gain \\
\midrule
None (full sample) & 22 & 7.536 & $+3.536$ \\
$z<0.08$ & 21 & 7.581 & $+3.581$ \\
$z<0.10$ & 21 & 7.581 & $+3.581$ \\
$z<0.12$ & 20 & 2.286 & $-1.714$ \\
$z<0.20$ & 19 & 0.871 & $-3.129$ \\
$z<0.30$ & 17 & 1.155 & $-2.845$ \\
\bottomrule
\end{tabular}
\end{table}

% ====================================================================
\subsection{Low-$z$ / Mid-$z$ Component Decomposition}
\label{sec:components}

To see which redshift zone drives the result, and to address the concern that the
low-$z$ arm overlaps the peculiar-velocity region, we fit the two Gaussian arms of
$T$ as independent nuisance amplitudes. Individually, the low-$z$ arm $G(z;0.09,0.06)$
gives $\Delta\chi^2=11.69$ and the mid-$z$ arm $G(z;0.775,0.10)$ gives
$\Delta\chi^2=4.81$; fitted together (two amplitudes) they give $\Delta\chi^2=17.13$,
essentially the single-template value, confirming the two arms are nearly independent
in the precision metric. The CPL-absorbing power is concentrated in the low-$z$ arm:
CPL adds only $\Delta\chi^2=0.88$ after the low-$z$ arm alone, but still
$\Delta\chi^2=7.45$ after the mid-$z$ arm alone---the mid-$z$ arm is nearly orthogonal
to CPL. Both arms nonetheless carry beyond-CPL structure ($5.04$ and $4.73$ added
after CPL, respectively).

Crucially, the cumulative node-removal test (Table~\ref{tab:cumulative}) localizes the
low-$z$ power to $z\approx0.10$--$0.20$, not to the lowest node $z=0.05$ where peculiar
velocities dominate: removing nodes only up to $z<0.10$ leaves CPL essentially
unchanged, and the preference collapses only once $z<0.12$--$0.20$ is removed. The
low-$z$ arm is thus tracing the $z\approx0.10$--$0.20$ survey-transition region---the
same low-$z$ range corrected by the Hoyt et al.~\cite{hoyt2026} host-mass re-measurement
(Section~6)---rather than the very-low-$z$ peculiar-velocity zone.

\begin{table}[h]
\centering
\caption{Decomposition of the template into its two Gaussian arms (fixed
$\Omega_m=0.2975$ protocol). The low-$z$ arm carries most of the CPL-absorbing power;
the mid-$z$ arm is nearly orthogonal to CPL.}
\label{tab:components}
\begin{tabular}{lccc}
\toprule
Component & $\Delta\chi^2$ alone & CPL after it & it after CPL \\
\midrule
Low-$z$ arm $G(z;0.09,0.06)$  & 11.69 & 0.88 & 5.04 \\
Mid-$z$ arm $G(z;0.775,0.10)$ & 4.81  & 7.45 & 4.73 \\
Both arms (2 amplitudes)      & 17.13 & --   & --   \\
Full template $T$             & 17.12 & 0.87 & 10.45 \\
\bottomrule
\end{tabular}
\end{table}

% ====================================================================
\subsection{Template--CPL Overlap (Degeneracy Test)}
\label{sec:overlap}

A natural concern is that the template $T(z)$ merely reproduces the CPL distance
deformation, so that its improvement is the dynamical--dark-energy signal relabeled.
We test this directly in the Union3 precision metric, with the marginalized offset $M$
projected out. The angle between $T$ and the CPL best-fit deformation
$\delta\mu_{\rm CPL}=\mu_{\rm CPL}-\mu_{\Lambda\rm CDM}$ is $\cos\theta=0.79$
($\theta\approx38^\circ$); projecting $T$ onto the full CPL tangent subspace
$\{\partial\mu/\partial w_0,\,\partial\mu/\partial w_a\}$ gives the same value, with
$63\%$ of $T$ lying inside the CPL subspace and $37\%$ orthogonal to it. The two
directions are thus strongly aligned but not degenerate. The nested $\Delta\chi^2$
decomposition (Table~\ref{tab:overlap}) is markedly asymmetric: once $T$ is included CPL
adds only $\Delta\chi^2=0.87$, whereas once CPL is included $T$ still adds
$\Delta\chi^2=10.45$. CPL is therefore almost entirely contained in $T$ ($88\%$ of its
pull is shared), while most of $T$'s improvement is structure CPL cannot produce,
localized at the $z\approx0.78$ survey transition (Figure~\ref{fig:overlap}).
Consistently, the recovered amplitude \emph{persists} when CPL is free in the model:
$a=0.0279$ for $\Lambda$CDM$+$nuisance (matching the compressed-data value) and
$a=0.0357$ for CPL$+$nuisance.

Were $T$ merely the dark-energy signal in disguise, $\theta$ would vanish, $T$ would add
nothing after CPL, and $a$ would collapse toward zero; none of these holds. This rules
out exact $T$--CPL degeneracy. It does not by itself distinguish a systematic from a
signal---goodness-of-fit cannot---which is why the discriminating evidence remains the
host-mass correspondence of Section~6 and the survey-composition
transitions, not the fit alone. The large $\Delta\chi^2=10.45$ that $T$ adds after CPL
reflects genuine residual structure rather than overfitting. The beyond-CPL null test of
Section~7.2 provides the direct overfitting check for this degeneracy concern: after CPL
is already fitted, the original template still adds $\Delta\chi^2=10.45$, exceeding every
random smooth template in the same after-CPL test (maximum $8.99$; add-one $p\approx1/201$).
The residual structure $T$ captures is therefore not typical of a random smooth direction
once the CPL-absorbable component has been removed. The quantity relevant to the
dark-energy claim is in any case the $0.87$ that CPL adds after $T$, not the $10.45$.

\begin{table}[h]
\centering
\caption{Template--CPL overlap (fixed $\Omega_m=0.2975$ protocol). The angle $\theta$
is measured in the Union3 precision metric with the magnitude offset $M$ marginalized.
All inputs reproduce the Table~\ref{tab:protocol} diagnostics to rounding precision.}
\label{tab:overlap}
\begin{tabular}{lc}
\toprule
Quantity & Value \\
\midrule
$\cos\theta$ ($T$ vs CPL deformation)          & $0.79$ \\
$\theta$ ($T$ vs CPL 2D subspace)              & $37.6^\circ$ \\
$T$ inside CPL subspace ($\cos^2\theta$)       & $63\%$ \\
$T$ orthogonal to CPL ($\sin^2\theta$)         & $37\%$ \\
CPL $\Delta\chi^2$ added after $T$             & $0.87$ \\
$T$ $\Delta\chi^2$ added after CPL             & $10.45$ \\
shared (overlap) $\Delta\chi^2$                & $6.67$ \\
amplitude $a$, $\Lambda$CDM$+$nuisance (mag)   & $0.0279$ \\
amplitude $a$, CPL$+$nuisance (mag)            & $0.0357$ \\
\bottomrule
\end{tabular}
\end{table}

\begin{figure}[h]
\centering
\figorbox{overlap_angle.png}{0.95\textwidth}{(a) T(z) vs its best CPL projection; (b) nested chi2 ledger: CPL after T = 0.87, T after CPL = 10.45.}
\caption{Template--CPL overlap in the Union3 precision metric.
\textbf{(a)} The nuisance template $T(z)$ (blue) and its best projection onto the CPL
subspace, $T_\parallel$ (red dashed). CPL reproduces $T$ at low redshift but cannot fit
the sharp $z\approx0.78$ step (shaded $T_\perp$, the $37\%$ of $T$ orthogonal to all
$w_0$--$w_a$ deformations). \textbf{(b)} Nested $\Delta\chi^2$ ledger: CPL adds only
$0.87$ after $T$, but $T$ adds $10.45$ after CPL: CPL is largely absorbed by $T$, while $T$ contains structure beyond CPL.}
\label{fig:overlap}
\end{figure}

% ====================================================================
\section{Independent Reproduction}

An independent reproduction of the Union3.1 UNITY1.8 results directly from the raw FITS
matrix (without astropy, using manual header parsing and big-endian float64 reading)
confirmed: $\Delta\chi^2(\mathrm{CPL})=4.7393$, $\Delta\chi^2(\text{Gaussian
template})=11.9581$, $\Delta\chi^2(\text{low-minus-mid step})=11.7335$, and
$\Delta\chi^2(\text{CPL after step})=0.1243$. All values match to rounding precision.
The reproduction also confirmed that node-removal tests must subselect the covariance
matrix $C=P^{-1}$ and re-invert; subsetting the precision matrix directly gives
incorrect results.

% ====================================================================
\section{Limitations}

\begin{itemize}
\item Compressed-data and public-data proxy analyses, plus a full-likelihood
cross-check in which the nuisance mode is added inside a public reimplementation of the
Union3 likelihood---not as an official, internally validated Union3/DESI systematic
model. These tests are not a replacement for the official DESI$+$CMB$+$Union3 MCMC
likelihood with an internally modelled systematic.
\item The full-likelihood cross-check uses Planck NPIPE CamSpec (not the official DESI
\texttt{plik} choice) and omits CMB lensing; the compressed cross-check represents CMB
by a Gaussian prior on $\Omega_m h^2$.
\item The Union3 compressed product distributes a posterior under a spline-node prior
\cite{kim2024}; this affects absolute contour placement more than relative model
comparisons.
\item The low-$z$ template component overlaps redshift ranges where peculiar velocities
and survey-composition changes can introduce coherent distance distortions; the
component decomposition (Section~\ref{sec:components}) localizes the CPL-absorbing power
to $z\approx0.10$--$0.20$ rather than the lowest-$z$ peculiar-velocity zone, but does not
fully exclude a low-$z$ systematic contribution.
\item These tests do not prove that DESI dynamical dark energy is false. They identify a
robustness concern in the Union3 contribution.
\item Analysis was conducted with AI assistance. All numerical results are reproducible
from public data and scripts.
\end{itemize}

% ====================================================================
\section{Recommended Falsification Test}

The definitive test is to rerun the official DESI $+$ CMB $+$ Union3 $w_0$--$w_a$
likelihood with an added low-$z$ / mid-$z$ redshift-dependent nuisance parameter and
compare $\Lambda$CDM$+$nuisance, CPL, and CPL$+$nuisance. If CPL loses AIC support after the nuisance is included, this would indicate that the Union3 contribution to the $w_0$--$w_a$ preference is strongly sensitive to this systematic-like direction, and would be consistent with a systematics-limited SN contribution. If the preference persists, the concern is resolved. The
full-likelihood cross-check of Section~\ref{sec:fullplanck} carries out this comparison
in the public pipeline and finds that CPL does lose AIC support
($+8.42\to-0.38$) once the nuisance is marginalized; an official internal-systematics
treatment is the remaining step. The three-version progression (Section~6) suggests
that as Union3.1 corrections are incorporated into the full joint likelihood, the CPL
preference may weaken further, consistent with the weakening already observed in
official Union3.1 joint analyses.

% ====================================================================
\section{Conclusion}

The Union3 data show sensitivity to a low-$z$ / mid-$z$ nuisance direction in CPL
robustness tests. A single Gaussian nuisance template improves the fixed-$\Omega_m$
residuals more than twice as well as CPL ($\Delta\chi^2=17.12$ vs $\approx7.52$) and
renders CPL AIC-disfavored once included. A compressed-prior Cobaya cross-check gives
the same qualitative picture (CPL $\Delta\chi^2$ $8.29\to0.52$), and a full-likelihood
cross-check using Planck NPIPE CamSpec, DESI DR1 BAO, and Union3 confirms it directly:
the strong AIC preference for CPL ($+8.42$ vs $\Lambda$CDM) is removed once the nuisance
is marginalized ($-0.38$ vs $\Lambda$CDM$+$nuisance), with the posterior shifting from
$2.4\sigma$ to $1.6\sigma$ from $\Lambda$CDM and a recovered amplitude
$a=0.029\pm0.010$ matching the compressed-data value. Template-grid, random-null, and
leave-one-region-out tests show that the effect is not driven by one exact functional
form or by one narrow redshift interval. Across Union3, Union3.1 UNITY1.7, and Union3.1
UNITY1.8, the nuisance amplitude decreases monotonically and the CPL AIC gain becomes
small, suggesting that the Union3 contribution to the DESI DR1 $w_0$--$w_a$ preference
is sensitive to a low-$z$ / mid-$z$ systematic direction that is reduced, but not
removed, in Union3.1. These results motivate a full likelihood reanalysis with an
explicit redshift-dependent nuisance parameter.

% ====================================================================
\begin{thebibliography}{9}
\small
\bibitem{desi2024} Adame et al.\ (DESI Collaboration) 2024,
\href{https://arxiv.org/abs/2404.03002}{arXiv:2404.03002} --- DESI DR1 BAO cosmological constraints.
\bibitem{rubin2025} Rubin D.\ et al.\ 2025,
\href{https://arxiv.org/abs/2311.12098}{arXiv:2311.12098}, ApJ --- Union3 $+$ UNITY1.5 supernova compilation.
\bibitem{hoyt2026} Hoyt T., Rubin D.\ et al.\ 2026,
\href{https://arxiv.org/abs/2601.19424}{arXiv:2601.19424} --- Union3.1 host-galaxy mass re-measurements and low-$z$ distance corrections.
\bibitem{rubin2026} Rubin D.\ et al.\ 2026,
\href{https://arxiv.org/abs/2601.19854}{arXiv:2601.19854} --- Union3.1 $+$ UNITY1.8 two-population light-curve standardization framework.
\bibitem{efstathiou2025} Efstathiou G.\ 2025, MNRAS, \href{https://arxiv.org/abs/2408.07175}{arXiv:2408.07175} --- Systematic offset in low-$z$ supernovae in Pantheon$+$ and DES5Y.
\bibitem{planck2018} Planck Collaboration 2018,
\href{https://arxiv.org/abs/1807.06209}{arXiv:1807.06209} --- Planck 2018 cosmological parameters.
\bibitem{kim2024} Kim A.G.\ 2024,
\href{https://arxiv.org/abs/2412.14181}{arXiv:2412.14181} --- Comments on the Union3 spline-interpolated distance-moduli model.
\bibitem{cortes2025} Cort\^es M.\ \& Liddle A.R.\ 2025,
\href{https://arxiv.org/abs/2504.15336}{arXiv:2504.15336} --- On DESI DR2 exclusion of $\Lambda$CDM; principled combination yields $3.1\sigma$.
\bibitem{data} Data: \texttt{rubind/union3\_release} (GitHub) --- Union3 and Union3.1 compressed likelihood files (public).
\bibitem{truong2026} Truong J., Aldering G., Perlmutter S., Rubin D., Schlegel D.\ 2026,
\href{https://arxiv.org/abs/2604.18859}{arXiv:2604.18859} --- Testing $\Lambda$CDM versus dynamical dark energy in one year: a DESI spectroscopic follow-up program for Rubin supernovae.
\bibitem{huang2025} Huang L., Cai R.-G., \& Wang S.-J.\ 2025,
\href{https://arxiv.org/abs/2502.04212}{arXiv:2502.04212}, Sci.\ China Phys.\ Mech.\ Astron.\ --- DESI dynamical-dark-energy evidence biased by low-$z$ supernovae.
\bibitem{ormondroyd2026} Ormondroyd A.N., Handley W.J., Hobson M.P., Lasenby A.N., \& Yallup D.\ 2026,
\href{https://arxiv.org/abs/2509.13220}{arXiv:2509.13220}, MNRAS --- Bayesian model selection between dark energy and supernova biases.
\bibitem{mirpoorian2026} Mirpoorian S.H., Lin M.-X., \& Pogosian L.\ 2026,
\href{https://arxiv.org/abs/2604.24761}{arXiv:2604.24761} --- Cosmological impact of redshift-dependent Type~Ia supernova calibration.
\bibitem{vincenzi2025} Vincenzi M.\ et al.\ (DES) 2025,
\href{https://arxiv.org/abs/2501.06664}{arXiv:2501.06664} --- Comparing DES-SN5YR and Pantheon$+$; offset largely from Malmquist-bias/intrinsic-scatter modelling.
\bibitem{popovic2025} Popovic B.\ et al.\ 2025,
\href{https://arxiv.org/abs/2511.07517}{arXiv:2511.07517} --- DES SN reanalysis (Dovekie recalibration); continued preference for evolving dark energy.
\end{thebibliography}

\vspace{1em}
\noindent\rule{\textwidth}{0.4pt}\\
\small Jaanus Hass $|$ Independent researcher, Saru, Estonia $|$ \texttt{jaanush@hotmail.com} $|$
May--June 2026 $|$ Reproduction scripts:
\href{https://github.com/JaanusHass/union3-lowz-midz-nuisance}{github.com/JaanusHass/union3-lowz-midz-nuisance}

\end{document}
